\documentclass[11pt,letterpaper]{article}
\usepackage[utf8]{inputenc}
\usepackage[T1]{fontenc}
\usepackage{lmodern}
\usepackage[margin=0.28in]{geometry}
\usepackage{hyperref}
\hypersetup{colorlinks=false, urlcolor=black, pdfborder=0 0 0}

\setlength{\parindent}{0pt}
\setlength{\parskip}{0pt}
\renewcommand{\baselinestretch}{0.82}
\setlength{\lineskip}{0pt}
\setlength{\lineskiplimit}{0pt}

\makeatletter
\renewcommand{\@listi}{%
  \leftmargin=\leftmargini
  \rightmargin=0pt
  \labelsep=4pt
  \labelwidth=16pt
  \itemindent=0pt
  \listparindent=0pt
  \topsep=0pt
  \partopsep=0pt
  \parsep=0pt
  \itemsep=-11pt
}
\let\saveditemize\itemize
\let\savedenditemize\enditemize
\renewcommand{\itemize}{\bgroup\setlength{\baselineskip}{0.65\baselineskip}\saveditemize}
\renewcommand{\enditemize}{\savedenditemize\egroup}
\makeatother

\begin{document}

{\Large \textbf{HSU, Ke-Cheng (Clio)}}\\[1pt]
(852) 6840 9753 \quad $|$ \quad \href{mailto:khsu@connect.ust.hk}{khsu@connect.ust.hk} \quad $|$ \quad \href{https://github.com/cliohsu1997}{github.com/cliohsu1997} \quad $|$ \quad \href{https://cliohsu1997.github.io/}{cliohsu1997.github.io}

\vspace{0.15pt}

{\large \textbf{Professional Summary}}\\[-0.5pt]
PhD candidate in Economics at HKUST (expected June 2027) and Hong Kong permanent resident, specializing in digital markets and consumer behavior. Analyze large proprietary datasets, including e-commerce platform transactions and retail point-of-sale records; collect data at scale through web scraping, APIs, and AI-based OCR in Python. Build reproducible, testable data pipelines on the cloud (AWS) and computing clusters (HPC/SLURM); present results to non-specialist audiences through clear writing, data visualizations, and weekly HTML reports. Apply demand and pricing models to what-if questions on customer experience, platform design, and growth.

\vspace{0.15pt}

{\large \textbf{Education}}\\[-0.5pt]
\textbf{The Hong Kong University of Science and Technology}, Hong Kong\\
\textit{Doctor of Philosophy in Economics} \hfill Expected Graduation: June 2027\\
Focus: how digital markets and retail decisions affect shoppers and sellers.\\
\textit{Master of Philosophy in Economics} \hfill 2020--2022\\
CGA: 4.1/4.3.\\
\textit{Bachelor of Science in Mathematics and Economics, Minor in Actuarial Mathematics} \hfill 2015--2019\\
CGA: 3.9/4.3; graduated first in class (Academic Achievement Medal, top 1\%).\\
Awards: University Scholarship (tuition and living allowance, 730K HKD, 2015--2019); Dean's List (2016--2019); Postgraduate Studentship (1.5M HKD); Red Bird PhD Award (40K HKD, 2022).

\vspace{0.15pt}

{\large \textbf{Experience}}\\[-0.5pt]
\textbf{Teaching Assistant \& Lecturer}, Department of Economics, HKUST \hfill 2020--present\\
\vspace{-6pt}
\begin{itemize}
  \item Lectured the MSc Math Camp (rated 4.53/5); led tutorials for PhD and undergraduate economics courses (average rating 4.72/5).
\end{itemize}
\textbf{Research Assistant} for Prof.\ Benjamin Bernard, National Taiwan University \hfill 2020\\
\vspace{-6pt}
\begin{itemize}
  \item Modeled a banking crisis as a negotiation between the regulator and banks' creditors; simulated when a promise not to bail out the banks is believable.
\end{itemize}

\vspace{0.15pt}

{\large \textbf{Projects}}\\[-0.5pt]
\textbf{Inventory and Pricing under Resale Price Maintenance (RPM)} \textit{(w/ K.\ Kawaguchi, HKUST; K.\ Kawai, U.\ Tokyo)}\\
\vspace{-6pt}
\begin{itemize}
  \item Cleaned and organized monthly bookstore point-of-sale records into a relational database: panels of deliveries, sales, and returns by title, market, and month (tens of billions of records across 1{,}000+ stores), linked via standardized ISBNs and merged with local demographics and title information scraped from Amazon in Python.
  \item Built a dynamic model of how bookstores price and stock titles under uncertain demand and RPM, capturing stockouts and unsold returns (35--40\% of Japanese books printed are discarded); estimated demand and supply with instrumental variables (2SLS), validated the estimates with Monte Carlo simulations, selected controls with LASSO, and implemented the model in C++.
  \item Showed through descriptive statistics and regressions that when titles stock out, shoppers switch to other titles or wait for restock rather than buying from other bookstores.
  \item Benchmarked the model against the standard static approach, which overstates seller costs by \(\sim\)20\% because it ignores demand uncertainty and inventory decisions.
  \item Distilled results into weekly HTML reports with charts and tables for collaborators.
\end{itemize}
\vspace{-3pt}
\textbf{Return Policy Design in E-Commerce} \textit{(Job Market Paper)}\\
\vspace{-6pt}
\begin{itemize}
  \item Built and estimated a model of the full customer journey --- purchase, return-shipping-insurance take-up, and return --- using retailer switches between automatic and selective insurance adoption on Taobao (2017--2023; \(\sim\)62M RMB revenue; six apparel categories; \(\sim\)20--30\% return rates).
  \item Developed a reproducible, version-controlled pipeline on AWS and HPC/SLURM (Python, R, C++): scraped product attributes and digitized product images with AI-based OCR; cleaned and linked orders, prices, insurance, and returns; estimated category-specific demand, consumer fit uncertainty, and return responses; accelerated simulations in C++; validated policy rankings (seller profit, consumer surplus, social welfare) with bootstrap tests.
  \item Identified a key trade-off: a stricter return policy curbs wasteful returns, but the resulting lower return rates make higher prices more attractive to sellers and can block valuable purchases; showed that providing return-shipping insurance to all shoppers works best in high-margin categories, letting each seller choose works better in low-margin ones (a \(\sim\)15--30 percentage-point gain), and a partial return penalty (discounted price plus \(\sim\)10--15\% restocking fee) reaches \(\sim\)75--78\% of the best-possible outcome outside dresses.
  \item Presented findings at the Asia-Pacific Industrial Organization Conference 2026, the Econometric Society Australasia Meeting 2026, and the China Economics Annual Conference Summer Forum 2026.
\end{itemize}

\vspace{0.15pt}

{\large \textbf{Skills}}\\[-0.5pt]
\textbf{Languages:} Mandarin (Native); Cantonese (Basic); English (Fluent); Taiwanese (Native).\\
\textbf{Technical:} Python, R, Stata, C++, Excel, Git; AWS, HPC/SLURM; relational databases; AI-assisted coding and testing.\\
\textbf{Communication:} Hall Tutor (conflict mediation and resident activities); Teaching Assistant Coordinator (coaching and evaluating a teaching team).\\
\textbf{References:} Kei Kawai (UTokyo; Compass Lexecon), Kohei Kawaguchi (HKUST), Pak Hung Au (HKUST).

\end{document}
